\documentclass[../../main.tex]{subfiles}
\begin{document}

{\bf [解]}
新しく
\begin{align}
 I_n=\displaystyle\int_0^\pi x^2|\sin nx|dx \\
 a_k=\displaystyle\int_{\frac{k-1}{n}\pi}^{\frac{k}{n}\pi}x^2|\sin nx|dx \\
 f(x)=x^2
\end{align}
として，$\left[\dfrac{k-1}{n}\pi,\dfrac{k}{n}\pi\right]$で$f(x)$の$\max,\min$を与える$x$を$M_k,m_k$とする．

$|\sin nx|\ge0$より
\begin{align}
f(m_k)|\sin nx|\le f(x)|\sin nx|\le f(M_k)|\sin nx|
\end{align}
だから，$\left[\dfrac{k-1}{n}\pi,\dfrac{k}{n}\pi\right]$の区間で積分すると
積分して，
\begin{align}
f(m_k)\int_{\frac{k-1}{n}\pi}^{\frac{k}{n}\pi}|\sin nx|dx\le a_k\le f(M_k)\int_{\frac{k-1}{n}\pi}^{\frac{k}{n}\pi}|\sin nx|dx \quad\cdots\text{①}
\end{align}
ここで，$\displaystyle\int_{\frac{k-1}{n}\pi}^{\frac{k}{n}\pi}|\sin nx|dx=\dfrac{2}{n}$だから代入して
\begin{align}
\dfrac{2}{n}f(m_k)\le a_k\le\dfrac{2}{n}f(M_k)
\end{align}
となる．$k$について和をとって，
\begin{align}
\dfrac{2}{n}\sum_{k=1}^nf(m_k)\le I_n\le\dfrac{2}{n}\sum_{k=1}^nf(M_k) \quad\cdots\text{②} \label{eq:2}
\end{align}
となる．左右について区分求積から極限値は
\begin{align}
\begin{split}
\dfrac{1}{n}\sum_{k=1}^nf(m_k)\longrightarrow\dfrac{1}{\pi}\int_0^\pi f(x)dx&=\dfrac{1}{3}\pi^2 \\
\dfrac{1}{n}\sum_{k=1}^nf(M_k)\longrightarrow\dfrac{1}{3}\pi^2
\end{split}
\end{align}
だから，\cref{eq:2}及びはさみうちの定理から，
\begin{align}
I_n\longrightarrow\dfrac{2}{3}\pi^2
\end{align}
である．


\newpage

\end{document}
